\chapter*{Phụ Lục I. Ba Mươi Lời Tự Răn}
\addcontentsline{toc}{chapter}{Phụ Lục I. Ba Mươi Lời Tự Răn}
\markboth{Ba Mươi Lời Tự Răn}{Ba Mươi Lời Tự Răn}

\begin{truyen}{Ba Mươi Lời Tự Răn Ngắn Gọn}{Thirty Short Self-Admonitions}
\chuhoa{N}{hững câu ngắn dưới đây không nhằm tạo cảm giác đạo lý nặng nề. Chúng được viết để đọc lại vào những lúc lòng bắt đầu mềm đi trước một khoản chi không cần thiết.}
\textit{(The short lines below are not meant to sound morally heavy. They are written to be reread whenever the heart begins to soften before an unnecessary expense.)}

1. Tôi không được tiêu tiền để mua sự quên lãng tạm thời.\
\textit{(I must not spend money to buy temporary forgetfulness.)}

2. Mỗi đồng tiền đi vào nhà đều có bóng dáng của lao động.\
\textit{(Every coin entering the home carries the shadow of labor.)}

3. Tiền nhỏ coi thường mãi sẽ thành nỗi lo lớn.\
\textit{(Small money repeatedly ignored becomes a large worry.)}

4. Không phải thích là nên mua.\
\textit{(Not everything liked should be purchased.)}

5. Mua sắm không chữa được mọi nỗi mệt mỏi.\
\textit{(Shopping does not heal every kind of exhaustion.)}

6. Sĩ diện thường đắt hơn giá món đồ.\
\textit{(Pride often costs more than the item itself.)}

7. Tôi không cần hơn người khác ở bề ngoài để gia đình mình yên hơn.\
\textit{(I do not need to look above others for my family to live in greater peace.)}

8. Tiền dạy thêm là mồ hôi ngoài giờ, không được tiêu như tiền nhặt được.\
\textit{(Tutoring money is after-hours sweat; it must not be spent like found money.)}

9. Không có quỹ dự phòng là tự đẩy mình đến gần hoảng loạn hơn.\
\textit{(Having no emergency fund pushes me closer to panic.)}

10. Người biết dừng tay đúng lúc thường đỡ phải gồng lưng về sau.\
\textit{(Those who stop their hand at the right moment usually strain less later.)}

11. Tiết kiệm không làm tôi nghèo niềm vui, nó giúp tôi giàu bình yên.\
\textit{(Saving does not make me poor in joy; it helps me become rich in peace.)}

12. Một lần bớt mua vô ích là một lần bảo vệ tương lai.\
\textit{(Each useless purchase avoided is an act of protecting the future.)}

13. Tôi không chờ còn dư mới tiết kiệm. Tôi phải tiết kiệm trước.\
\textit{(I do not wait for leftovers to save. I must save first.)}

14. Nợ chỉ nên đi cùng tính toán, không nên đi cùng cảm xúc.\
\textit{(Debt should accompany calculation, not emotion.)}

15. Cái thiếu lớn nhất đôi khi không phải tiền, mà là kỷ luật.\
\textit{(The greatest shortage is sometimes not money, but discipline.)}

16. Con tôi học cách dùng tiền từ dáng sống của tôi.\
\textit{(My children learn how to use money from the way I live.)}

17. Một cuốn sổ chi tiêu thành thật quý hơn nhiều lời hứa đẹp.\
\textit{(An honest expense notebook is worth more than many beautiful promises.)}

18. Tôi không được lấy tiền tương lai vá cho sự yếu lòng hôm nay.\
\textit{(I must not use future money to patch today’s weakness.)}

19. Không phải ai có nhiều đồ hơn cũng sống vững hơn.\
\textit{(Not everyone who owns more things lives more steadily.)}

20. Biết đủ không làm tôi nhỏ lại, nó làm lòng tôi yên lại.\
\textit{(Knowing enough does not shrink me; it calms my heart.)}

21. Một căn nhà gọn gàng và vừa sức đáng quý hơn một vẻ ngoài hào nhoáng.\
\textit{(A home that is modest and orderly is worth more than outward glamour.)}

22. Nếu món đồ này phải đổi bằng nhiều giờ dạy mệt mỏi, tôi cần nghĩ lại.\
\textit{(If this item costs many exhausting teaching hours, I need to think again.)}

23. Tôi không cần tự thưởng cho mình bằng mọi ham muốn nảy ra.\
\textit{(I do not need to reward myself through every desire that appears.)}

24. Mỗi lần tiêu theo cảm xúc là một lần mở cửa cho hối tiếc.\
\textit{(Every emotional purchase opens the door to regret.)}

25. Người lớn không chỉ biết kiếm tiền; người lớn còn biết giữ tiền.\
\textit{(An adult does not only know how to earn money; an adult also knows how to keep it.)}

26. Gia đình cần sự ổn định hơn những món đồ gây ấn tượng.\
\textit{(A family needs stability more than impressive possessions.)}

27. Tiền để dành không nằm im vô ích; nó đang đứng gác.\
\textit{(Saved money is not lying idle; it is standing guard.)}

28. Tôi chỉ có thể dạy người khác điều mình đang tập sống.\
\textit{(I can teach others only what I am trying to live.)}

29. Không ai cứu nổi tôi khỏi lãng phí nếu chính tôi không chịu dừng lại.\
\textit{(No one can rescue me from waste if I refuse to stop myself.)}

30. Muốn ngày mai đỡ khổ, hôm nay phải bớt dễ dãi với chính mình.\
\textit{(If tomorrow is to suffer less, today I must be less indulgent with myself.)}
\end{truyen}

\begin{baihoc}
Những câu tự răn ngắn có tác dụng khi được đọc lặp lại vào đúng lúc lòng bắt đầu dao động.
\textit{(Short self-admonitions work when they are reread at the exact moment the heart begins to waver.)}
\end{baihoc}

\begin{ghinhoanh}
Repeat the right sentence often.

Language can train the hand.
\end{ghinhoanh}

\ngancach